\section{Un modèle statistique de l'écart de portage Totem--FO}\label{sec-noise}

\subsection{Modèle d'observation}
\begin{definition}[Modèle d'observation]\label{def-obs}
Le détecteur observe l'écart de portage $y\in\R^K$ de \cref{eq-y} et le modélise par
\begin{equation}\label{eq-obs}
y=X\beta+\sum_{t\in\mathcal{S}}a_t\,t+\varepsilon,\qquad X=\bigl[\,T,\;T\ln T\,\bigr],\qquad\varepsilon\sim\N(0,\kappa\Sigma),
\end{equation}
où $X\beta$ est un écart de repo lisse $\delta b(T)=\beta_0+\beta_1\ln T$ écrit dans l'espace de portage (désaccord
légitime, y compris une véritable erreur de marquage FO), $\mathcal{S}$ un ensemble éventuellement vide de gabarits
d'anomalie (\cref{sec-detection}), $\Sigma$ la covariance du bruit de consensus construite ci-dessous à partir des
statistiques Totem, et $\kappa>0$ une échelle unique calibrée sur un historique sain.
\end{definition}
$\beta_0$ est l'écart de repo à un an ($\ln1=0$) ; $\beta_1$ sa pente en $\ln T$. Sous l'hypothèse nulle,
$\mathcal{S}=\emptyset$.

\subsection{Covariance des moyennes de consensus}
D'après le \cref{lem-additivity} et la \cref{prop-gap}, le bruit de consensus en portage vaut
$\varepsilon_k=-e^F_k-e^D_k$, où $e^F_k$ est l'erreur du consensus du panel options $\ln\overline{F}_k$ et
$e^D_k=\delta\DB(T_k)/(S-D_k)$ celle du consensus dividendes alloué. Les deux panels étant indépendants
(\cref{ass-panels} avec $m=0$, le cas prudent),
\begin{equation}\label{eq-sigma}
\Sigma=\Sigma^F+\Sigma^D.
\end{equation}
Les deux termes sont des covariances de moyennes de panels emboîtés, donc la \cref{prop-nested} s'applique : chacun est
une covariance au niveau des contributeurs divisée terme à terme par $\max(n_k,n_l)$. Il reste à modéliser la
covariance au niveau des contributeurs à partir des dispersions publiées par Totem.

\subsection{Le panel options}
L'erreur de log-forward d'un contributeur en $T$ vaut $-T\,\delta b_j-\delta D_j(T)/(S-D)$. Les vues de repo sont lisses
(les desks marquent quelques piliers et interpolent) : la partie repo est essentiellement un niveau persistant, de
covariance $\sigma_b^2T_kT_l$. La partie dividendes est cumulative : un facteur persistant (un contributeur optimiste sur
les dividendes l'est pour toutes les années) plus des incréments qui arrivent avec chaque dividende.

\begin{lemma}[Covariance cumulative]\label{lem-cumcov}
Pour $s\in\R^K_{\ge0}$ croissant et $w\in[0,1]$, la matrice $C_{kl}=w\,s_ks_l+(1-w)\min(s_k,s_l)^2$ est semi-définie
positive : c'est la covariance de $Z_k=\sqrt{w}\,s_k\,\xi+\sqrt{1-w}\,W(s_k^2)$, avec $\xi\sim\N(0,1)$ indépendant d'un
mouvement brownien standard $W$. Pour $w=0$, la corrélation entre $k<l$ vaut $s_k/s_l$ (incréments indépendants).
\end{lemma}
\begin{proof}
$\Cov(Z_k,Z_l)=w\,s_ks_l+(1-w)\Cov\bigl(W(s_k^2),W(s_l^2)\bigr)=w\,s_ks_l+(1-w)\min(s_k^2,s_l^2)$, et une matrice de
covariance est semi-définie positive. Avec $w=0$, $\Cov(Z_k,Z_l)/\sqrt{\Var Z_k\Var Z_l}=s_k^2/(s_ks_l)=s_k/s_l$ pour $k<l$.
\end{proof}

\begin{proposition}[Covariance du panel options]\label{prop-sigmaF}
Soit $s_k=\sigma_{F,k}/F_k$ la dispersion publiée des forwards des contributeurs, en logarithme, et $\sigma_b$ la
dispersion des vues de repo. En décomposant $s_k^2=\sigma_b^2T_k^2+s_k'^2$,
\begin{equation}\label{eq-sigmaF}
\Sigma^F_{kl}=\frac{\sigma_b^2\,T_kT_l+w\,s'_ks'_l+(1-w)\min(s'_k,s'_l)^2}{\max(n_{A,k},n_{A,l})}.
\end{equation}
\end{proposition}
\begin{proof}
L'indépendance des vues de repo et de dividendes donne la somme d'un terme de rang un et de la covariance cumulative du
\cref{lem-cumcov} pour la partie dividendes ; la \cref{prop-nested} donne la division par le plus grand panel.
\end{proof}

\begin{remark}[Estimation de $\sigma_b$]
Aux échéances les plus courtes, l'incertitude sur les dividendes est négligeable, donc $s_k\approx\sigma_bT_k$. Le
prototype prend la médiane de $s_k/T_k$ sur les trois premiers piliers. Le minimum serait biaisé vers le bas (minimum de
trois estimations bruitées) : il sous-estimerait l'incertitude sur le niveau de repo et ferait déclencher trop souvent le
test d'erreur de marquage FO de la \cref{sec-detection}. Ce défaut a été constaté sur les titres difficiles à emprunter,
puis corrigé.
\end{remark}

\subsection{Le panel dividendes}
\begin{proposition}[Covariance du panel dividendes]\label{prop-sigmaD}
Supposons que les erreurs de consensus par période aient la covariance
$(\Sigma_L)_{yz}=\rho_{yz}\,\sigma_{D,y}\sigma_{D,z}/\max(n_{B,y},n_{B,z})$, avec $\rho_{yz}=e^{-|y-z|/\ell}$. Alors, au premier
ordre,
\begin{equation}\label{eq-sigmaD}
\Sigma^D=J_A\,\Sigma_L\,J_A^\top,\qquad (J_A)_{ky}=\frac{A_{ky}}{S-D_k}.
\end{equation}
\end{proposition}
\begin{proof}
$e^D_k=\delta D_k/(S-D_k)=\sum_yA_{ky}\delta L_y/(S-D_k)$ d'après le \cref{lem-allocation} : une application linéaire des
erreurs par période.
\end{proof}

\subsection{Les pics de dispersion sont des symptômes, pas du bruit}
\begin{proposition}[Masquage]\label{prop-masking}
Dans le désaccord de date de la \cref{prop-timing}, soit $s_k$ la dispersion hors désaccord et $\Delta=\tilde d_i/(S-D_k)$.
La statistique de test du déplacement en $T_k$, calculée avec la dispersion publiée (gonflée), vaut
\begin{equation}\label{eq-masking}
z=\frac{p\,\Delta}{\sqrt{s_k^2+p(1-p)\Delta^2}/\sqrt{n_{A,k}}},
\end{equation}
qui reste bornée par $\sqrt{n_{A,k}\,p/(1-p)}$ quelle que soit la taille du dividende, alors qu'avec la dispersion hors
désaccord $z$ croît linéairement en $\Delta$.
\end{proposition}
\begin{proof}
La \cref{prop-timing}(i) donne le déplacement moyen et (iii) la variance gonflée. Quand $\Delta\to\infty$, le rapport
tend vers $\sqrt{n_{A,k}}\,p/\sqrt{p(1-p)}$.
\end{proof}

\begin{example}
Pour l'exemple de la galerie ($p=0{,}37$, $n_A=16$), la borne vaut $\sqrt{16\times0{,}37/0{,}63}\approx3{,}1$ : à peine
le seuil critique de la famille de tests, si bien que le désaccord passerait souvent inaperçu. Le prototype plafonne donc
un pic isolé de dispersion à $1{,}5$ fois la plus grande des deux dispersions voisines avant de construire $\Sigma$, et
publie le rapport, qui est lui-même un indice de désaccord entre contributeurs, en général sur une date de détachement
(ici $9{,}7$).
\end{example}

\begin{insight}
$\Sigma$ n'utilise que ce que Totem publie (dispersions et nombres de contributeurs) plus trois réglages structurels :
le poids du facteur $w$, la longueur de corrélation $\ell$ entre périodes de dividendes et la dispersion du repo lue sur
le court terme. L'échelle unique $\kappa$ absorbe un biais commun de ces réglages ; c'est leur \emph{forme} qui rend le
détecteur puissant, car elle lui indique quels écarts sont bon marché (lisses, longs) et lesquels sont coûteux (courts,
locaux).
\end{insight}
